\problemname{KIT Finding}

\illustration{0.3}{fox.jpg}{
    CC BY-SA 4.0 by Uoaei1 on \href{https://commons.wikimedia.org/wiki/File:Vulpes_vulpes_Mallnitz_01.jpg}{Wikimedia Commons}%
    \vspace{-5mm}% Make sure the first sample fits on the first page.
}

``\emph{Find the Fox}'' is a recent book that contains $200$ pages of word search puzzles comprising only the letters `\texttt{F}', `\texttt{O}' and `\texttt{X}'.
The special feature of the book is that there is only a single occurrence of the word ``\texttt{FOX}'' throughout all of its pages.

Recall that in a word search the goal is to find hidden words (only one hidden word in this case) in a grid of letters.
The words can occur horizontally or vertically or diagonally, as well as forwards or backwards, allowing for a total of $8$ different reading directions.

For this year's NWERC, we want to create a baby version of ``\emph{Find the Fox}''.
Your goal in this problem, therefore, is to create a word search grid with given dimensions and containing each of the letters `\texttt{K}', `\texttt{I}' and `\texttt{T}' a given number of times.
Similar to the original book, this grid should contain exactly one occurrence of the word ``\texttt{KIT}''.

\begin{figure}[!h]
  \centering
  \includegraphics[width=0.4\textwidth]{sample.pdf}
  \caption{Illustration of Sample Output 1. Flip page for the solution.}
\end{figure}

\begin{Input}
The input consists of:
\begin{itemize}
\item One line with five integers $h$, $w$, $k$, $i$, and $t$ ($3 \le h,w \le 100$, $k,i,t \ge 1$, $k+i+t = h\cdot w$), where $h$ and $w$ are the height and width the word search grid should have, and $k$, $i$, and $t$ specify the required number of occurrences of `\texttt{K}', `\texttt{I}' and `\texttt{T}', respectively.
\end{itemize}
\end{Input}

\begin{Output}
  Output a word search grid according to the given rules.
  It can be shown that such a word search grid always exists.

  If there are multiple valid solutions, you may output any one of them.
\end{Output}

\vfill
\begin{figure}[!h]
  \centering
  
\includegraphics[angle=180,width=0.25\textwidth]{sample-solved.pdf}
  \caption{Illustration of the second sample.}
\end{figure}
\vspace{2cm}
